% Short variant of sec_intro_task.tex (S1, S1.1, S2). Material removed here is
% relocated verbatim to Appendix C; nothing is dropped.
\section{Introduction}

Ask a language model to place 21 circles in a unit square so as to maximize the sum of their
radii, and it will not search. It reaches for the nearest square grid --- five by five, radius
$1/(2k)$ --- then truncates it to 21 circles, leaving four cells empty and their area unclaimed.
A better construction is one parameter away: drop to a four-by-four grid and add corner fillers of
radius $(\sqrt{2}-1)/(2k)$, and the sum rises. The model does neither, and more often than not the
next sample repeats the same move.

The behavior can be written down. Writing $k^*(N) = \mathrm{round}(\sqrt{N})$ for the nearest-square
order, a value function $V(k, m)$ over grid order $k$ and filler count $m$ identifies the modal
output of the proposal distribution --- which template the model emits most often at each $N$, and
hence its value, from the problem parameters alone, before sampling. The empirical mode equals the
predicted value at all seven $N$ tested, and per-sample agreement is bounded by the modal frequency
itself, 56--86\% by cell (\S3.2). The rule also predicts where the behavior costs value: when
$k^{*2} \le N$ the model extends the grid with fillers (a branch arm \S3.4 later restricts to
$m \le 1$); when $k^{*2} > N$ it truncates instead of dropping to $k^*{-}1$ and filling. These
\emph{trap zones} --- $N \in [13,15]$, $[21,24]$, $[31,35]$, $[43,48]$, $[57,63]$ (clipped to
$[57,60]$ by our sweep bound) --- are a property of the branch behavior, not of the problem.

Circle packing is the showcase benchmark of the LLM-driven discovery literature, the lineage
FunSearch \citep{romeraparedes2024funsearch} opened. It is the task on
which AlphaEvolve \citep{novikov2025alphaevolve} reported 2.63586276 for $n = 26$ and
on which later systems report a narrow band --- ShinkaEvolve \citep{lange2025shinkaevolveopenendedsampleefficientprogram} 2.6359831 (its relaxed-tolerance figure; 2.6359777 strict), HELIX
2.63598308 \citep{su2026helixevolutionaryreinforcementlearning}, GigaEvo 2.636 \citep{khrulkov2025gigaevoopensourceoptimization}, AdaEvolve 2.636 \citep{cemri2026adaevolveadaptivellmdriven} --- with
ThetaEvolve \citep{wang2025thetaevolvetesttimelearningopen} claiming new best-known bounds rather than matching it.

\textbf{What we measure, and what we do not.} Those systems query $p(\text{program} \mid \text{parent programs} +
\text{fitness scores} + \text{evaluator feedback} + \text{archive context, at generation } t > 0)$. We query
$p(\text{coordinates} \mid \text{task description, no parent, no code channel, no feedback, generation 0})$. These
differ on four axes at once --- parent conditioning, output modality, feedback, generation index
--- and we do not claim to have measured the in-loop distribution. A source audit finds no cited
system making the unconditioned call: every in-loop call in FunSearch, ShinkaEvolve and OpenEvolve
is program-conditioned, and the nearest deployed relatives are trivial-seed initializations plus
AlphaEvolve's ``No evolution'' ablation (Appendix~C.1). The unconditioned call is the limit these
approach as seed information goes to zero; that this limit is a template lookup with a computable
value is actionable for loop builders --- seed diversification, forced-$k$ initialization, trap-$N$
avoidance in benchmark selection, and a generation zero on which the downstream diversity machinery
\citep{mouret2015mapelites,lehman2011novelty} is doing repair rather than exploration
(Appendix~C.1). Anything about what the loop samples at $t > 0$ is a hypothesis here. The one
published result bearing directly on the substitution --- ``Dictionaries, Not Darwin''
\citep{li2026dictionariesdarwinsetlevelselection}, reporting parent-conditioned evolution
indistinguishable from fresh independent sampling at equal budget --- is warrant from equation
discovery, not geometry. \S8 names the experiment that would close the gap.

Three framing choices are worth stating plainly. Behavioral anchoring is established territory
\citep{huang2026understandinganchoringeffectllm,lou2024anchoringbiaslargelanguage,malberg2025comprehensiveevaluationcognitivebiases}
and we inherit its vocabulary; what differs is modality --- the anchor is a \emph{construction
template}, not a number. The benchmark choice is strategic, not a novelty claim. Preregistration,
too, is an adopted standard here
\citep{thomas2026mitigatingllmbasedphackingpreregistering,williams2026predictingllmsafetyrelease,vaccaro2026preregistrationexperimentsaiagents,zhang2026testingfrontierlargelanguage};
our twist is what is locked --- exact-output point predictions from a closed form, on a held-out
container.

\textbf{Contributions.}
\begin{enumerate}
\item \emph{A choice/execution dissociation: the model selects the better construction and cannot build it through either channel.} Handed the provably higher-valued in-family construction with its score, the weak tier chooses it, then fails to instantiate it in 30 of 31 invalid attempts, the remaining one unevaluable (arm CH, \S3.5 --- a disclosed post-hoc attempt-level reading whose classifier is the family filler radius; the registered per-cell prediction P-CH2 held at one of three cells); delegated to an executed math-only program, the same ceiling holds across two tiers and a fresh-draw replication --- 0 of 115 valid program outputs exceed the family argmax (arms CC/CC2/CCS, \S3.7). The template is not what the model prefers but what it can reliably execute, in text or in code. This finding does not depend on whether any deployed loop runs unconditioned calls.
\item \emph{A closed form that names the modal model output in advance.} $k^*(N)$ with $V(k, m)$ and
$T(k, N)$ (the extend-with-fillers and truncated-grid values, \S2.2) identify the value a weak-tier
proposer emits most often at a given $N$, verified against a linear-programming oracle on 83
configurations to within $10^{-9}$ and tested out of sample on two containers, the rectangle rule
restated ($q^* = \mathrm{round}(\sqrt{N/a})$, $p^* = \mathrm{round}(\sqrt{N \cdot a})$) but never
refitted. Prior work publishes functional forms predicting \emph{aggregate} metrics ahead of a run
\citep{kaplan2020scalinglawsneurallanguage,hoffmann2022trainingcomputeoptimallargelanguage,snell2024predictingemergentcapabilitiesfinetuning};
to our knowledge none predicts a specific multi-decimal \emph{individual emitted output value}
ahead of sampling. The claim is scoped to the weak tier (\S5).
\item \emph{A tier boundary condition, as three attractor families.} Ambition rises monotonically
with nominal tier; validity does not. Matched cells, primary tolerance: 83\% $\rightarrow$ 100\%
$\rightarrow$ 13\% (78\% over all seven weak-tier cells). At $10^{-9}$: 64\% $\rightarrow$ 90\%
$\rightarrow$ 13\%, the 64\% over all seven weak-tier cells, its matched-cell counterpart not
computed (Table~\ref{tab:ladder}). The most ambitious tier attempts recursive gaskets and mostly
fails to produce a valid packing; it was addressed through a bare serving alias and appears as
\texttt{opus\_alias} with that caveat everywhere.
\item \emph{Checkable faithfulness, and trace requests as interventions.} Method lines are
checkable against emitted coordinates, and 54 of 56 scoreable claims (96.4\%) describe the object
actually built (Wilson 95\% interval $[88\%, 99\%]$, \S6.5). Requesting the line at all is an
intervention rather than an observation: the concentration shift it produced (87\% vs 70\%,
$p = 0.0325$ uncorrected) fails multiplicity correction, is carried by one of three cells and is
wave-confounded (\S6.4), so it is reported as met-as-registered exploratory material, not a
demonstrated effect. Studies collecting process descriptors \emph{by request} are nonetheless
measuring a perturbed distribution.
\item \emph{A family boundary, measured from both sides --- and a near-deterministic strong form.}
Anchoring is not a single-vendor artifact: it transfers to two further vendors' weak models
(61--67\% of valid outputs against 56--76\% in-family at the same five cells --- overlapping
ranges at these $n$, \S3.2 and \S3.6), while one further family brackets it --- gemma-4-26b escapes \emph{upward} to the
in-family rival construction (27/43 valid outputs at discriminating cells, those where prediction
and family argmax differ) and gemma-4-31b, above the same validity floor, anchors on nothing
(\S3.6). Where anchoring holds, its strong form is near-deterministic: pooling every weak-tier
discriminating-cell validity across the original-family and transfer-vendor arms, the provably
better in-family rival is emitted 3 times in 186 (arm F 2/57, trace 1/53, rectangle 0/11, arm M
0/10, arm V 0/16, arm CN 0/39 --- a descriptive pool over separately registered weak-tier arms,
each reported at its own bar in \S3--\S4 and \S6; the Sonnet tier reaches the rival's structure 5
times in 30 and is excluded by tier, \S5; gemma-4-26b, which does \emph{not} anchor at its
discriminating cells, is the counterexample that proves the statistic is a family trait, \S3.6).
\end{enumerate}

\textbf{Non-claim guard.} Everything here is a behavioral regularity over emitted outputs. We make
no claim about mechanism inside the weights, and the paper uses that vocabulary throughout: the
model \emph{emits}, the distribution \emph{concentrates}, the formula \emph{identifies the modal
output} (storage-and-retrieval vocabulary corrected). \S8 states the leading mechanism-free alternative, arithmetic tractability, and its direct
probe (arm CH, \S3.5): wrong about selection, right about instantiation.

\subsection{Trap zones}

Substituting $k^*(N)$ into the branch condition: TRAP $N \in [k^2-k+1, k^2-1]$, CONVERGE
$N \in [k^2, k^2+k]$. Over $N = 10\ldots60$: $[13,15]$, $[21,24]$, $[31,35]$, $[43,48]$ and
$[57,63]$ clipped to $[57,60]$ by the sweep bound. Inside a zone the rule generally loses value
against the best construction available \emph{within the recipe family}. Worst-in-zone penalties
fall as $N$ grows: 8.51\% at $N = 13$, 7.03\% at 21, 6.01\% at 31, 5.25\% at 43, 4.66\% at 57. The
penalty reaches zero at $N = 35$, 48, 14 and 15, but \textbf{not} at $N = 24$, where truncation
gives $T(5,24) = 2.4000000$ against $V(4,8) = 2.4142136$, a residual 0.59\%. Branch label and
penalty are therefore not coextensive, and 4 of the 22 trap-branch $N$ in range cost nothing.
Figure~\ref{fig:trapzones} plots the prediction against the family optimum and the published bounds
over the whole sweep; it and the rest of the trap-zone arithmetic --- the independent
linear-programming verification of every closed-form value over 83 configurations to within
$10^{-9}$, and the published-bound comparison with its $N > 30$ coverage gap --- are in
Appendix~C.2. \S3.3 uses $N = 35$ as the zero-penalty control.

\section{Task and recipe family}

\subsection{The benchmark}

Maximize the sum of radii of $N$ non-overlapping circles in the unit square. Feasibility is
decidable exactly from emitted coordinates, so proposals score with no human in the loop. All
invocations in the core arms are zero-shot and code-free: the prompt (Appendix A) fixes the count,
states containment and non-overlap, forbids writing or executing code, and demands only a raw
Python list of \texttt{[x, y, r]} triples. Our reason for the code-free design was that a code
channel routes the task to an executed optimizer, so the distribution would reflect the optimizer
--- a reason now measured. Three preregistered code-channel arms (\S3.7) show that a math-only
executed program neither clears the family ceiling (0 of 115 valid outputs above the argmax, both
tiers) nor unseats the anchor as the modal weak-tier output, so the code-free scope costs less loop
relevance than assumed. The library-assisted channel (\texttt{numpy}/\texttt{scipy}), the one
deployed loops permit, stays unmeasured.

\subsection{The recipe family}

Three terms, used consistently: the \textbf{recipe family} is the parametric set of grid-plus-filler
constructions; a \textbf{template} is one member; the \textbf{anchor} is the template the
distribution concentrates on at a given $N$. A $k \times k$ grid places one circle per cell center
with $r_{\mathrm{grid}} = 1/(2k)$; $m$ fillers sit on interior grid vertices, tangent to their four
surrounding grid circles, with $r_{\mathrm{filler}} = (\sqrt{2} - 1)/(2k)$, $0 \le m \le (k-1)^2$.
Hence $V(k, m) = k/2 + m(\sqrt{2} - 1)/(2k)$. When $N < k^2$ the observed behavior is not to drop to
a smaller grid and fill it but to \emph{truncate} --- lay out the $k \times k$ lattice, occupy $N$
cells, leave the radius unchanged --- giving $T(k, N) = N/(2k)$. These reproduce every previously
reported anchor with no fitting: $V(5,1) = 2.5414214$, $V(5,2) = 2.5828427$, $V(5,0) = 2.500$,
$T(5,23) = 2.300$, $V(4,7) = 2.3624369$. The prior-arm and ledger evidence that the family is where
the weak tier lives --- including the two sub-$10^{-6}$ exceedances among 155 rows --- is in
Appendix~C.3.

\subsection{The order rule is a definition; the branch rule is the hypothesis}

We write $k^*(N) = \lfloor \sqrt{N} + \tfrac{1}{2} \rfloor = \mathrm{round}(\sqrt{N})$ for the
nearest-square order. This is a \emph{definition}, not a fitted law: two formalizations of
``nearest'' agree on the integers, and the discrete selection among candidate orders is only weakly
identified (Appendix~C.4). The falsifiable content is the \textbf{branch rule}:
$k^{*2} \le N \rightarrow$ extend with $m = N - k^{*2}$ fillers (converge);
$k^{*2} > N \rightarrow$ truncate the $k^*$-grid (trap). The extend arm survived testing only at
$m \le 1$ --- a preregistered extension disconfirmed it at $m \ge 4$, where the model prefers the
$(k^*{+}1)$ rectangle grid (\S3.4). Nothing about ``nearest'' forces truncation --- a proposer with
any local search would drop to $k^*{-}1$ and fill, which is why the trap zones exist and why a
higher-valued rival is reachable inside them. The governing quantity is the signed distance
$N - k^{*2}$; primality is not the variable --- 13, 23, 31, 43, 47 and 59 trap while 11, 17, 19,
29, 37, 41 and 53 converge --- refuting the prior work's conjecture about primes near 30 before any
model is queried.

\subsection{Notation and conventions}

Terms used throughout, gathered here so no verdict table depends on a definition made elsewhere:

\begin{center}
\begin{footnotesize}
\begin{tabular}{@{}p{0.24\linewidth}p{0.70\linewidth}@{}}
\toprule
Term & Meaning \\
\midrule
$k^*(N)$ & nearest-square grid order, $\mathrm{round}(\sqrt{N})$ (a definition, \S2.3) \\
$V(k, m)$, $T(k, N)$ & extend-with-$m$-fillers value $k/2 + m(\sqrt{2}-1)/(2k)$; truncated-grid value $N/(2k)$ \\
anchor & the template the proposal distribution concentrates on at a given $N$ \\
attractor & the same object, used when the emphasis is on the family a tier converges to rather than on one cell (\S5) \\
on-prediction & emitted sum within $2\times10^{-3}$ of the registered predicted value \\
rival (argmax) & the highest-valued member of the recipe family at that $N$ when it differs from the prediction --- inside a trap zone, $V(k^*{-}1, m)$ \\
discriminating cell & $N$ where prediction $\ne$ family argmax; only these can distinguish anchoring from family search \\
validity & non-overlap and containment at $10^{-6}$ (primary) and $10^{-9}$ (logged), both always reported \\
trap / converge & branch by sign of $N - k^{*2}$: truncate when $k^{*2} > N$, extend when $k^{*2} \le N$ \\
tier & nominal capability class of the serving alias addressed (weak/mid/top); \texttt{opus\_alias} is a bare alias with no attestable weights binding \\
BELOW-FLOOR & model failing its arm's registered validity floor; no anchoring claim in either direction \\
UNSCOREABLE & cell with $<$ 3 valid samples under the GM-chain registration; claims nothing \\
mixed-serving accounting & pooled rate including cells collected through a registered serving-path amendment; the first-party-only rate is always given beside it \\
\bottomrule
\end{tabular}
\end{footnotesize}
\end{center}
